\bigskip
\bigskip
\subsubsection*{Խնդիրներ և հարցեր թեմա 13-ի վերաբերյալ}

\bigskip
\bigskip
\subsubsection*{Խնդիրներ և հարցեր թեմա 13-ի վերաբերյալ}

\begin{enumerate}[label=\thesection.\arabic*.]

% 13.1
\item Ապացուցեք $(\mathbb{R}, \text{սովոր.})$ թվային ուղղի $(-\infty; 0] \cup (0; +\infty)$ տրոհումից ստացվող ֆակտոր-տարածությունը հոմեոմորֆ է կապակցված երկկետանի (Սիրպինսկու) տարածությանը (սահմանումը տես թեմա 4-ում):

% 13.2
\item Դիտարկենք $(\mathbb{R}, \text{սովոր.})$ տարածության կետերի միջև $R$ էկվիվալենտության հարաբերությունն, ըստ որի $x_1 R x_2 \Leftrightarrow (x_1 - x_2)$-ը ամբողջ թիվ է: Ապացուցեք՝
\begin{enumerate}
    \item[ա)] $R$-ը համարժեքության հարաբերություն է;
    \item[բ)] $\mathbb{R}/R$ ֆակտոր-տարածությունը հոմեոմորֆ է $S^1$ շրջանագծին:
\end{enumerate}

\begin{hint}
Օգտվել թեմա 2-ում օրինակ 9-ից և այս թեմայի օրինակ 2-ից:
\end{hint}

% 13.3
\item Դիտարկենք $\mathbb{R}^2$ կոորդինատային հարթությունը սովորական տոպոլոգիայով և $R$ էկվիվալենտության հարաբերություն նրա կետերի միջև, ըստ որի $(x_1, y_1) R (x_2, y_2) \Leftrightarrow x_1 - x_2$ և $y_1 - y_2$ թվերը ամբողջ են: Ապացուցեք՝
\begin{enumerate}
    \item[ա)] $R$-ը համարժեքության հարաբերություն է;
    \item[բ)] $\mathbb{R}^2/R$ ֆակտոր-տարածությունը հոմեոմորֆ է տորին:
\end{enumerate}

\begin{hint}
Օգտվել թեմա 2-ում օրինակ 10-ից:
\end{hint}

% 13.4
\item Ճի՞շտ է արդյոք պնդումը. եթե ունենք տոպոլոգիական տարածությունների $f: X \to Y$ անընդհատ, սյուրյեկտիվ արտապատկերում, ապա $Y$-ը $X$-ի ֆակտոր-տարածությունն է (որոշված $f$ արտապատկերմամբ):
\par Լուծում: Դիտարկենք $(X, \text{դիսկրետ})$, $(Y, \text{անտիդիսկրետ})$ տարածությունները, որտեղ $X=\{x_1, x_2\}, Y=\{y_1, y_2\}$ և $f: X \to Y, f(x_1)=y_1, f(x_2)=y_2$ արտապատկերումը: Ցույց տվեք, որ $Y$ տոպոլոգիան չի համընկնում $f$ արտապատկերմամբ որոշվող $Y$-ի ֆակտոր-տոպոլոգիայի հետ:

% 13.5
\item Ապացուցեք, որ $\mathbb{RP}^1$ պրոյեկտիվ ուղիղը հոմեոմորֆ է շրջանագծին:

\begin{hint}
Ներկայացրեք $\mathbb{RP}^1$-ը որպես փակ կիսաշրջանագծի ֆակտոր-տարածություն:
\end{hint}

% 13.6
\item Ստուգեք, արդյոք բաց (փակ) արտապատկերում է օրինակ 1-ում $P: X \to X/R$ կանոնական արտապատկերումը:

% 13.7
\item Ճի՞շտ է արդյոք պնդումը. ցանկացած $X/R$ ֆակտոր-տարածության դեպքում $P: X \to X/R$ կանոնական արտապատկերումը բաց արտապատկերում է:
\par Լուծում: Դիտարկենք օրինակ 2-ում $[0, 1/2) \subset [0, 1]$ ենթատարածության կերպարը կանոնական արտապատկերման դեպքում:

% 13.8
\item Ապացուցեք. կամայական $X/R$ ֆակտոր-տարածության դեպքում $P: X \to X/R$ կանոնական արտապատկերումը բաց (փակ) արտապատկերում է այն և միայն այն դեպքում, երբ $X$-ի ցանկացած $U$ բաց ենթաբազմության դեպքում $P^{-1}(P(U))$ ենթաբազմությունը բաց է (փակ է) $X$-ում:

\par Մինչև հաջորդ խնդիրներին անցնելը տանք մի սահմանում:
\par Դիցուք $X$-ը տոպոլոգիական տարածություն է, իսկ $A$-ն $X$-ի որևէ ենթաբազմություն է: Դիտարկենք $R$ համարժեքության հարաբերություն $X$-ի վրա համարելով՝
\begin{enumerate}
    \item[ա)] տեղի ունի $a_1 R a_2$ հարաբերություն $A$-ի ցանկացած երկու $a_1, a_2$ տարրերի դեպքում;
    \item[բ)] եթե $x_1, x_2 \in X \setminus A$, կամ $x_1 \in A$ և $x_2 \in X \setminus A$, ապա տեղի ունի $x_1 R x_2$ միայն և միայն այն դեպքում, երբ $x_1 = x_2$: Այսպիսով համարժեքության դասերը $X$-ի մի կետանոց ենթաբազմություններն են և $A$-ն:
\end{enumerate}
\par Այս դեպքում $X/R$ ֆակտոր-տարածության մասին ասում են, որ այն ստացվում է $X$-ից սեղմելով մի կետի $A$-ի բոլոր կետերը: Նշանակվում են այդ տարածությունը $X/A$ սիմվոլով:

\begin{enumerate}[label=\thesection.\arabic*., resume]

% 13.9
\item Դիցուք $X$-ը $\mathbb{R}^2$ կոորդինատային հարթությունում $B^2 = \{(x, y) | x^2+y^2 \le 1\}$ միավոր շրջանն է, իսկ $A$-ն նրա եզրային $S^1 = \{(x, y) | x^2+y^2 = 1\}$ միավոր շրջանագիծն է: Ապացուցեք, որ $X/A$ ֆակտոր-տարածությունը հոմեոմորֆ է $S^2 = \{(x, y, z) | x^2+y^2+z^2=1\}$ սֆերային:
\par \red{Պատկերավոր ասած՝ պարկը կապում ենք, ստացվում է սֆերա:}

\begin{hint}
Նախ կառուցեք հոմեոմորֆիզմ $B^2 = \{(x, y) | x^2+y^2 < 1\}$ ամբողջ շրջանի և առանց մի կետ $S^2 \setminus \{(0, 0, 1)\}$ շրջանի միջև:
\end{hint}

% 13.10
\item Դիցուք $X$-ը $x^2+y^2=1$ ուղիղ շրջանային գլանն է $\mathbb{R}^3$-ում, իսկ $A$-ն $x^2+y^2=1, z=0$ շրջանագիծն է: Ապացուցեք, որ $X/A$ ֆակտոր-տարածությունը հոմեոմորֆ է $Y = \{(x, y, z) | x^2+y^2-z^2=0\}$ կոնական մակերևույթին:
\par Պատկերավոր ասած՝ գլանում մի շրջանագիծ կետի սեղմելիս ստացվում է կոնական մակերևույթ:
\par Լուծում: Նախ ցույց տվեք, որ $(x, y, z) \mapsto (xz, yz, z)$ համադրումը որոշում է հոմեոմորֆիզմ $X \setminus A$ և $Y \setminus \{(0, 0, 0)\}$ տարածությունների միջև:

\end{enumerate}
